\subsection{直和的外代数。分次模的外代数}

设 A 为一个交换环， \(M = \bigoplus_{\lambda \in L} M_{\lambda}\) 为一族 A-模的直和，以及 \(j_{\lambda}: M_{\lambda} \to M\) 为典范单射；由此导出一个分次代数的 A-同态 \(\bigwedge(j_{\lambda}): \bigwedge(M_{\lambda}) \to \bigwedge(M)\) 。由于 \(\bigwedge(M)\) 是反交换的，§4 第 7 号命题 10（或在必要时推广到 L 为无限的情形，参见 §4 第 8 号注记 1）可应用于同态 \(\bigwedge(j_{\lambda})\) ；则存在唯一的代数同态

\[
g \colon \bigotimes_ {\lambda \in L} ^ {g} \Lambda (M _ {\lambda}) \to \Lambda (M)\tag{14}
\]

（也记作 \(g_{\mathbf{M}}\) ）使得 \(\bigwedge (j_{\lambda}) = g \circ f_{\lambda}\) ，其中

\[
f _ {\lambda}: \bigwedge (\mathrm{M} _ {\lambda}) \rightarrow \bigotimes_ {\lambda \in L} ^ {g} \bigwedge (\mathrm{M} _ {\lambda})
\]

表示典范同态。

命题10. 典范同态 \(g\) （公式（14））是分次代数同构。

为了证明 \(g\) 是双射，我们定义一个分次代数同态

\[
h: \bigwedge (\mathbf {M}) \rightarrow \bigotimes_ {\lambda \in L} ^ {\mathfrak {g}} \bigwedge (\mathbf {M} _ {\lambda})\tag{15}
\]

使得 \(g \circ h\) 与 \(h \circ g\) 分别是 \(\bigwedge(\mathbf{M})\) 与 \(\bigotimes_{\lambda \in L} \bigwedge(\mathbf{M}_{\lambda})\) 上的恒等映射。对于每个 \(\lambda \in L\) ，考虑复合线性映射

\[
u _ {\lambda}: \mathrm{M} _ {\lambda} \xrightarrow {\phi_ {\mathrm{M} _ {\lambda}} ^ {n}} \bigwedge (\mathrm{M} _ {\lambda}) \xrightarrow {f _ {\lambda}} \mathbf {g} _ {\lambda \in L} \bigotimes \bigwedge (\mathrm{M} _ {\lambda}).
\]

存在且仅存在一个 A-线性映射 \(u: \mathbf{M} \to {}^{\mathfrak{g}} \bigotimes_{\lambda \in L} \bigwedge (\mathbf{M}_{\lambda})\) ，使得 \(u \circ j_{\lambda} = u_{\lambda}\) （对所有 \(\lambda \in L\) ）。斜张量积 \({}^{\mathfrak{g}} \bigotimes_{\lambda \in L} \bigwedge (\mathbf{M}_{\lambda})\) 是一个交错代数：当 L 有限时，这由 §4 第 9 号命题 14 得出；当 L 任意时，这由 §4 第 8 号注记 1 中给出的该乘积的定义以及交错分次代数的正向极限是交错的这一事实得出。由于 \(u(\mathbf{M})\) 也包含在分次代数 \(\bigotimes_{\lambda \in L} \bigwedge (\mathbf{M}_{\lambda})\) 中次数为1的元素子模中，由第1号命题1和注记1可知，存在唯一的分次代数同态。

\[
h: \bigwedge (\mathbf {M}) \rightarrow \bigotimes_ {\lambda \in L} ^ {\mathfrak {g}} \bigwedge (\mathbf {M} _ {\lambda})
\]

使得 \(h \circ \phi_{\mathbf{M}}'' = u\) 。 \(g \circ h\) 与 \(h \circ g\) 为恒等映射这一事实的验证，与 §6 第 6 号命题 9 中的做法相同，只需将 \(\mathbf{S}\) 换为 \(\bigwedge\) ，将 \(\phi'\) 换为 \(\phi''\) 。

注记。设 \(N = \bigoplus_{\lambda \in L} N_{\lambda}\) 为以 L 为指标集的另一族 A-模的直和，并且对所有 \(\lambda \in L\) ，设 \(v_{\lambda}: M_{\lambda} \to N_{\lambda}\) 为一个 A-线性映射，于是存在 A-线性映射 \(v = \bigoplus_{\lambda} v_{\lambda}: M \to N\) （II，§1，第 6 号，命题 7）。则图表

\[
\begin{array}{c c c} \mathbf {g} _ {\lambda \in L} ^ {\otimes} \wedge (\mathrm{M} _ {\lambda}) & \xrightarrow {g _ {\mathrm{M}}} & \wedge (\mathrm{M}) \\ \bigotimes_ {\lambda \in L} \wedge (v _ {\lambda}) \Big \downarrow & & \Big \downarrow \wedge (v) \\ \mathbf {g} _ {\lambda \in L} ^ {\otimes} \wedge (\mathrm{N} _ {\lambda}) & \xrightarrow {g _ {\mathrm{N}}} & \wedge (\mathrm{N}) \end{array}
\]

是交换的（参见第4节第5号命题8的推论）。

\(\mathbf{g}_{\lambda \in L}^{\otimes} \bigwedge (\mathbf{M}_{\lambda})\) 的 A-子模（ \(\bigwedge^{n}(\mathbf{M})\) 借助同构 \(g\) 与之等同者）可以得到更精确的描述。对每个有限子集 \(J\) （ \(L\) 的），我们记 \(\mathrm{E}_{J} = \mathbf{g}_{\lambda \in J}^{\otimes} \bigwedge (\mathbf{M}_{\lambda})\) ，从而 \(\mathbf{g}_{\lambda \in L}^{\otimes} \bigwedge (\mathbf{M}_{\lambda}) = \lim_{\lambda \to 0} \mathrm{E}_{J}\) 相对于有向集 \(\mathfrak{F}(L)\) （ \(L\) 的有限子集组成的）成立，这是依据定义（§ 4，第 8 号，注记 1）。对每个族 \(\nu = (n_{\lambda}) \in \mathbf{N}^{(L)}\) （从而具有有限支集），它满足 \(\sum_{\lambda \in L} n_{\lambda} = n\) ；对每个有限子集 \(J\) （ \(L\) 的、包含族 \(\nu\) 的支集者），我们记

\[
\bigwedge^ {J, v} (\mathbf {M}) = \bigotimes_ {\lambda \in J} \bigwedge^ {n _ {\lambda}} (\mathbf {M} _ {\lambda})\tag{16}
\]

使得子模 \(\mathbf{E}_{\mathbf{J},n}\) ——即次数 \(n\) 的元素在 \(\mathbf{E}_{\mathbf{J}}\) 中组成的子模——是诸 \(\bigwedge^{\mathbf{J},\nu}(\mathbf{M})\) 的直和，这里族 \(\nu\) 的支集包含于 \(\mathbf{J}\) 且使得

\[
\sum_ {\lambda \in L} n _ {\lambda} = n
\]

（§ 4，第 7 号，命题 10 及第 8 号）。按约定，我们记 \(\bigwedge^{J,\nu}(\mathbf{M}) = \{0\}\) （对族 \(\nu\) ，其支集不包含于 \(J\) ）；于是

\(E_{J,n}\) 也可称为所有 \(\bigwedge^{J,\nu}(M)\) 的直和，其中 \(\nu\) 遍历集合 \(H_n\) ，即所有族 \(\nu = (n_\lambda)_{\lambda \in L}\) （满足 \(\sum_{\lambda \in L} n_\lambda = n\) 者）组成的集合。由于 \(\bigwedge^0(M_\lambda)\) 与 A 等同，显然对 L 的两个有限子集 \(J \subset J'\) 以及支集包含于 J 的族 \(\nu\) ，典范映射 \(\bigwedge^{J,\nu}(M) \to \bigwedge^{J',\nu}(M)\) （即典范映射 \(E_J \to E_{J'}\) 在 \(\bigwedge^{J,\nu}(M)\) 上的限制）是双射。如果对所有 \(\nu \in H_n\) ，

\[
\bigwedge^ {\nu} (\mathbf {M}) = \varinjlim \bigwedge^ {J, \nu} (\mathbf {M})\tag{17}
\]

于是可见，考虑到 II，§ 6，第 2 号，命题 5，我们有：

推论. A-模 \(\bigwedge^{n}(\mathbf{M})\) 是 \(\bigotimes_{\lambda \in L} \bigwedge (\mathbf{M}_{\lambda})\) 的一个子模在同构 (14) 下的像，该子模是诸子模 \(\bigwedge^{\nu}(\mathbf{M})\) 的直和，其中族 \(\nu = (n_{\lambda}) \in \mathbf{N}^{(L)}\) 满足 \(\sum_{\lambda \in L} n_{\lambda} = n\) ；若 \(J\) 是 \(\nu, \bigwedge^{\nu}(\mathbf{M})\) 的支集，则它典范同构于 \(\bigotimes_{\lambda \in J} \bigwedge^{n_{\lambda}}(\mathbf{M}_{\lambda})\) .

一般将 \(\bigwedge^{\nu}(\mathbf{M}),\bigotimes_{\lambda \in J}\bigwedge^{n_{\lambda}}(\mathbf{M}_{\lambda})\) 与它们在 \(\bigwedge^n (\mathbf{M})\) 中的像等同起来。按此约定：

命题11. 设 \(\Delta\) 为交换幺半群， \(\mathbf{M}\) 为类型 \(\Delta\) 的分次 A-模， \((\mathbf{M}_{\alpha})_{\alpha \in \Delta}\) 为其分次。对于每个有序对 \((\alpha, n) \in \Delta \times \mathbf{N}\) ，设 \(\bigwedge^{\alpha, n}(\mathbf{M})\) 为 \(\bigwedge^n (\mathbf{M})\) 的这样的子模：它是诸子模 \(\bigotimes_{\lambda \in J} \bigwedge^{n_{\lambda}}(\mathbf{M}_{\alpha_{\lambda}})\) 的直和，其中 \((n_{\lambda})_{\lambda \in L}\) 遍历整数 \(\geqslant 0\) 的、满足 \(\sum_{\lambda \in L} n_{\lambda} = n\) 的族所成的集合， \(J\) 为其支集，且对于每个 \((n_{\lambda})\) , \((\alpha_{\lambda})_{\lambda \in J}\) 遍历 \(\Delta^J\) 中满足 \(\sum_{\lambda \in J} \alpha_{\lambda} = \alpha\) 的族所成的集合。则 \((\bigwedge^{\alpha, n}(\mathbf{M}))_{(\alpha, n) \in \Delta \times N}\) 是唯一的类型为 \(\Delta \times N\) 、与 \(\bigwedge(\mathbf{M})\) 上的代数结构相容并且在 \(\mathbf{M} = \bigwedge^1 (\mathbf{M})\) 上诱导出给定分次的分次。

\(\bigwedge (\mathbf{M})\) 是诸 \(\bigwedge^{\alpha ,n}(\mathbf{M})\) 的直和这一事实由命题 10 的推论得出；证明的其余部分与 §5 第 5 号命题 7 的证明结尾完全相同。

